\documentclass[12pt]{article}

\usepackage[spanish, es-noshorthands]{babel}
\usepackage[a4paper, margin=3cm]{geometry}
\usepackage[T1]{fontenc}
\usepackage[utf8]{inputenc}
\usepackage{xcolor}
\usepackage{amsmath}
\usepackage{graphicx}
\usepackage{enumitem}
\usepackage{hyperref}
\usepackage{listings}
\usepackage{caption}

\definecolor{codebg}{RGB}{248,248,248}
\definecolor{codeframe}{RGB}{210,210,210}
\definecolor{codekeyword}{RGB}{0,92,150}
\definecolor{codestring}{RGB}{160,60,0}
\definecolor{codecomment}{RGB}{95,95,95}
\definecolor{codeident}{RGB}{70,70,70}

\lstdefinestyle{reportstyle}{
    backgroundcolor=\color{codebg},
    frame=single,
    rulecolor=\color{codeframe},
    basicstyle=\ttfamily\small,
    keywordstyle=\color{codekeyword}\bfseries,
    stringstyle=\color{codestring},
    commentstyle=\color{codecomment}\itshape,
    identifierstyle=\color{codeident},
    showstringspaces=false,
    numbers=left,
    numberstyle=\tiny\color{gray},
    stepnumber=1,
    numbersep=8pt,
    tabsize=4,
    breaklines=true,
    breakatwhitespace=true,
    columns=fullflexible,
    captionpos=b,
    keepspaces=true,
    xleftmargin=1em,
    framexleftmargin=0.8em
}

\lstset{style=reportstyle}

\hypersetup{
    colorlinks=true,
    linkcolor=codekeyword,
    urlcolor=codekeyword,
    pdftitle={Reporte de compilador de mini-lenguaje},
    pdfauthor={Ortega Novoa Octavio}
}

\begin{document}

\begin{titlepage}
\begin{center}

\begin{minipage}{0.49\textwidth}
\begin{flushleft}
\includegraphics[scale=0.12]{Imagenes/UNAM_logo.png}
\end{flushleft}
\end{minipage}
\begin{minipage}{0.49\textwidth}
\begin{flushright}
\includegraphics[scale=0.18]{Imagenes/LogoIngenieria.png}
\end{flushright}
\end{minipage}

\vfill

{\scshape\Large \textbf{UNIVERSIDAD NACIONAL AUTÓNOMA DE MÉXICO}} \\[4mm]
{\scshape\large \textbf{FACULTAD DE INGENIERÍA}} \\[8mm]

{\Large \textbf{Compiladores}} \\[4mm]
{\large Semestre 2026-2} \\[11mm]

{\Huge \textbf{Proyecto: Compilador de un mini-lenguaje con funciones, estructuras de control y expresiones booleanas}} \\[20mm]

\begin{tabular}{rl}
\textbf{Profesor:} & MANUEL ENRIQUE CASTAÑEDA CASTAÑEDA \\[2mm]
\textbf{Alumnos:} & Ortega Novoa Octavio \\
                  & Bailón Cruz Owen Isaac \\
                  & Torres Chávez Juan Carlos \\[4mm]
\textbf{Fecha de entrega:} & 29 de Mayo de 2026 \\
\end{tabular}

\vfill
\end{center}
\end{titlepage}

\tableofcontents
\newpage

\section{Introducción}
Este reporte presenta el diseño e implementación de un compilador para un mini-lenguaje de programación con fines académicos. Durante el desarrollo se buscó que el compilador no solo reconociera programas válidos, sino que también mostrara con claridad cómo una entrada fuente se transforma en código intermedio y después en un código final más simple.

La idea principal fue construir un traductor dirigido por la sintaxis capaz de reconocer declaraciones, expresiones aritméticas y booleanas, estructuras de control avanzadas (incluyendo selección múltiple), funciones sin parámetros, declaración de estructuras agrupadas y rutinas de entrada y salida estándar con soporte para cadenas de texto. 

El proyecto se fue armando por etapas: primero el análisis léxico, después el análisis sintáctico, luego la generación de código intermedio, más tarde la optimización y, al final, la traducción a un pseudocódigo de ensamblador simple. Una parte importante del trabajo fue ir probando cada fase por separado mediante una batería de scripts en C. Esto permitió detectar errores más rápido, corregirlos sin rehacer todo el compilador y entender mejor cómo se conectan las distintas partes del proceso de compilación.

Para llevarlo a cabo se usaron C como lenguaje principal de implementación, Flex para el análisis léxico, Bison para el análisis sintáctico y LaTeX para la documentación del reporte. Con ello se pudo construir tanto el compilador como el presente informe técnico.

\section{Requisitos del proyecto}
El proyecto contempla los siguientes componentes:
\begin{itemize}[label={-}]
    \item Declaraciones y uso de funciones sin parámetros (con y sin valor de retorno).
    \item Expresiones aritméticas y booleanas con operadores básicos.
    \item Estructuras de control de flujo: \texttt{if-else}, \texttt{while}, \texttt{do-while}, \texttt{for} y \texttt{switch/case}.
    \item Declaraciones y asignaciones de variables enteras, flotantes, de carácter y booleanas.
    \item Agrupación de datos mediante la palabra reservada \texttt{struct}.
    \item Entrada y salida estándar con lectura de variables e impresión de variables y cadenas literales.
    \item Generación de tabla de símbolos detallada.
    \item Generación de código intermedio de tres direcciones y código final tipo ensamblador.
    \item Optimización básica mediante propagación de constantes y eliminación de código muerto.
\end{itemize}

\section{Diseño del mini-lenguaje}
Antes de implementar el compilador se definió una sintaxis sencilla, fuertemente inspirada en C, para facilitar la escritura de programas de prueba y su traducción. Se eligió este estilo porque resulta familiar y evita introducir sobrecarga cognitiva con gramáticas exóticas mientras se trabaja en el núcleo del proyecto: la traducción semántica.

La intención no fue inventar un lenguaje completamente nuevo, sino diseñar uno lo bastante claro para poder probar funciones, condiciones, ciclos, agrupaciones y expresiones lógicas complejas de manera robusta.

\subsection{Elementos sintácticos principales}
El lenguaje incluye:
\begin{itemize}[label={-}]
    \item Tipos primitivos: \texttt{int}, \texttt{float}, \texttt{char}, \texttt{bool} y \texttt{void}.
    \item Operadores aritméticos: \texttt{+}, \texttt{-}, \texttt{*}, \texttt{/}, \texttt{++}, \texttt{--}.
    \item Operadores lógicos y relacionales: \texttt{\&\&}, \texttt{||}, \texttt{!}, \texttt{<}, \texttt{>}, \texttt{<=}, \texttt{>=}, \texttt{==}.
    \item Estructuras de control cíclico y condicional: \texttt{if}, \texttt{else}, \texttt{while}, \texttt{do-while}, \texttt{for}.
    \item Estructuras de selección múltiple: \texttt{switch}, \texttt{case}, \texttt{default}, \texttt{break}.
    \item Instrucciones de E/S: \texttt{read} y \texttt{print} (con soporte para mensajes entre comillas).
    \item Declaraciones de estructuras de datos con \texttt{struct}.
\end{itemize}

\subsection{Ejemplo de programa}
\begin{lstlisting}[language=C, caption={Ejemplo representativo del mini-lenguaje}]
struct Punto
{
    int x;
    int y;
};

int main()
{
    int a = 2;
    int b = 3;
    int c = a + b;

    print("Iniciando validacion");

    switch(c)
    {
        case 5:
            print("Resultado esperado");
            break;
        default:
            print("Resultado inesperado");
            break;
    }

    return c;
}
\end{lstlisting}

\section{Análisis léxico}
El análisis léxico se implementó utilizando Flex. Su función es reconocer las unidades léxicas (lexemas) del código fuente, asignarles un tipo de token y entregarlos al analizador sintáctico. Esta etapa es la primera capa de validación del proyecto y es la responsable de generar tablas de salida útiles para la depuración, como la tabla de tokens y la tabla de clasificación de variables y cadenas.

\subsection{Objetivo del lexer}
El lexer identifica:
\begin{itemize}[label={-}]
    \item Identificadores de variables, funciones y estructuras.
    \item Literales numéricos (enteros y flotantes), booleanos (\texttt{true}/\texttt{false}) y cadenas de texto.
    \item Palabras reservadas del lenguaje.
    \item Operadores matemáticos, lógicos y relacionales.
    \item Delimitadores y símbolos especiales (llaves, paréntesis, puntos y coma).
\end{itemize}

\subsection{Fragmento representativo del lexer}
\begin{lstlisting}[language=C, caption={Reconocimiento de tokens principales en Flex}]
"int"       { registerToken(yytext, "Palabras Reservadas"); return INT; }
"switch"    { registerToken(yytext, "Palabras Reservadas"); return SWITCH; }
"struct"    { registerToken(yytext, "Palabras Reservadas"); return STRUCT; }
"print"     { registerToken(yytext, "Palabras Reservadas"); return PRINT; }

{letter}+(_|{letter}|{digit})* { 
    registerToken(yytext, "Variables"); 
    yylval.string_val = strdup(yytext); 
    return IDENTIFIER;
}

\"([^\"\n]|\\.)*\" { 
    registerToken(yytext, "Cadenas"); 
    yylval.string_val = strdup(yytext); 
    return CHAR_LITERAL; 
}

"&&"        { registerToken(yytext, "Operadores Logicos"); yylval.string_val = strdup(yytext); return AND; }
\end{lstlisting}

\subsection{Comentarios de implementación}
Una decisión arquitectónica fundamental fue adjuntar el valor textual del lexema a múltiples tokens mediante \texttt{yylval.string\_val}. Esto permite que el parser no solo valide la gramática, sino que utilice directamente el valor de identificadores, literales y cadenas capturadas para construir la representación del código intermedio. Se descartaron los espacios en blanco y los comentarios en línea (\texttt{//}) para evitar ruido durante el análisis sintáctico.

\section{Análisis sintáctico}
El parser se implementó con Bison. Su función es validar la estructura jerárquica del programa y construir, de manera simultánea, los atributos y fragmentos de código necesarios mediante la técnica de traducción dirigida por la sintaxis. Aquí no solo importa que el programa esté bien escrito, sino que las acciones semánticas embebidas transformen cada instrucción correctamente.

\subsection{Estructura general}
La gramática define reglas robustas para:
\begin{itemize}[label={-}]
    \item El programa principal como una secuencia recursiva de sentencias.
    \item Declaraciones de variables, estructuras (\texttt{struct\_statement}) y funciones.
    \item Control de flujo complejo, destacando la integración de \texttt{switch\_statement} y sus respectivos \texttt{case}.
    \item Expresiones aritméticas y booleanas respetando la precedencia tradicional de operadores.
\end{itemize}

\subsection{Precedencia y asociatividad}
Para evitar la ambigüedad inherente de las gramáticas (como el clásico problema del \texttt{else} ambiguo), se declararon precedencias explícitas:
\begin{itemize}[label={-}]
    \item \texttt{OR} (\texttt{||}) y \texttt{AND} (\texttt{\&\&}) para operaciones lógicas.
    \item Comparaciones relacionales (\texttt{<}, \texttt{>}, \texttt{==}).
    \item Suma (\texttt{+}) y resta (\texttt{-}).
    \item Multiplicación (\texttt{*}) y división (\texttt{/}).
    \item \texttt{NOT} (\texttt{!}) como operador unario de alta prioridad.
\end{itemize}

\subsection{Fragmento del parser}
\begin{lstlisting}[language=C, caption={Reglas principales de sentencias en Bison}]
statement:
    declaration SEMICOLON
    | function_declaration
    | RETURN operation SEMICOLON
    | PRINT OPEN_PARENTHESIS operation CLOSE_PARENTHESIS SEMICOLON
    | READ OPEN_PARENTHESIS IDENTIFIER CLOSE_PARENTHESIS SEMICOLON
    | if_statement
    | while_statement
    | do_while_statement SEMICOLON
    | for_statement
    | switch_statement
    | struct_statement SEMICOLON
    | BREAK SEMICOLON
    ;
\end{lstlisting}

\subsection{Traducción dirigida por la sintaxis}
Cada regla sintáctica construye atributos sintetizados (\texttt{code}, \texttt{place}, \texttt{initCode}, etc.) que permiten ensamblar el código intermedio a medida que el árbol de derivación se resuelve. Esta decisión práctica elimina la necesidad de realizar pasadas adicionales sobre un Árbol Sintáctico Abstracto (AST) completo en memoria, produciendo variables temporales, etiquetas y saltos condicionales en tiempo real.

\section{Implementación y resolución de problemas}
El proyecto no se construyó de una sola vez; su desarrollo fue iterativo y por capas. Primero se validó el análisis léxico y sintáctico con programas triviales. Posteriormente, se incorporó la tabla de símbolos para controlar el alcance y la unicidad de las declaraciones. Una vez estabilizado este núcleo, se implementó la generación de código intermedio, la optimización y la traducción final.

Durante este proceso se resolvieron diversos desafíos técnicos:
\begin{itemize}[label={-}]
    \item \textbf{Propagación de lexemas:} Asegurar que el lexer entregara correctamente las cadenas literales con comillas y los nombres de variables sin truncarlos.
    \item \textbf{Ambigüedades gramaticales:} Solucionar conflictos \textit{shift/reduce} en expresiones booleanas combinadas y bloques condicionales anidados.
    \item \textbf{Redundancias de código:} Evitar la generación de etiquetas huérfanas o saltos consecutivos inútiles generados por construcciones como \texttt{do-while} o \texttt{switch}.
\end{itemize}

\section{Tabla de símbolos}
La tabla de símbolos es la estructura central que almacena la información semántica del programa: nombre, tipo y categoría (variable, función o estructura) de cada identificador. Su papel es fundamental para prevenir colisiones de nombres y validar que los símbolos se utilicen de acuerdo con su naturaleza.

\subsection{Ejemplo de uso interno}
\begin{lstlisting}[language=C, caption={Función de registro de símbolos}]
void addSymbol(char* name, char* type, char* category) 
{
    for(int i = 0; i < symbolCount; i++) 
    {
        if(strcmp(symbolTable[i].name, name) == 0) 
        {
            printf("[Error Semantico] La variable/funcion '%s' ya fue declarada.\n", name);
            return;
        }
    }

    strcpy(symbolTable[symbolCount].name, name);
    strcpy(symbolTable[symbolCount].type, type);
    strcpy(symbolTable[symbolCount].category, category);
    symbolCount++;
}
\end{lstlisting}

\section{Código intermedio de tres direcciones}
La generación de código intermedio se diseñó en formato de tres direcciones. Esta representación desvincula la lógica abstracta del lenguaje fuente de las limitaciones de una arquitectura de hardware específica, facilitando drásticamente los pasos de optimización.

\subsection{Características de la salida}
\begin{itemize}[label={-}]
    \item Uso exhaustivo de variables temporales (\texttt{t1}, \texttt{t2}, etc.) para desglosar expresiones matemáticas complejas.
    \item Generación de etiquetas autoincrementables (\texttt{L1}, \texttt{L2}, etc.) para el control de saltos.
    \item Transcripción directa de cadenas de texto para las instrucciones \texttt{print}.
    \item Instrucciones atómicas del tipo: \texttt{ifFalse tX goto LY}.
\end{itemize}

\section{Optimización básica}
La fase de optimización opera analizando textualmente el código intermedio generado para aplicar reducciones locales sencillas, garantizando un código más limpio antes de la traducción final.

\begin{itemize}[label={-}]
    \item \textbf{Propagación de constantes:} Si una variable se inicializa con un valor constante puro, ese valor reemplaza a la variable en operaciones subsecuentes dentro del bloque.
    \item \textbf{Plegado de constantes (Constant Folding):} Operaciones aritméticas entre literales (ej. \texttt{t1 = 2 + 3}) se resuelven en tiempo de compilación (\texttt{t1 = 5}).
    \item \textbf{Eliminación de código muerto:} El optimizador detecta instrucciones inalcanzables que suceden inmediatamente después de directivas \texttt{return} o \texttt{break} incondicionales, descartándolas del archivo de salida.
\end{itemize}

\section{Código final}
El código final toma el intermedio optimizado y lo mapea a un pseudocódigo de ensamblador didáctico. No busca acoplarse a una arquitectura física estricta (como x86 o ARM), sino demostrar la linealización de la memoria y la ejecución atómica de un procesador hipotético.

\subsection{Mapeo de mnemónicos}
\begin{itemize}[label={-}]
    \item Transferencia de datos: \texttt{MOV}.
    \item Aritmética: \texttt{ADD}, \texttt{SUB}, \texttt{MUL}, \texttt{DIV}.
    \item Lógica y comparación: \texttt{AND}, \texttt{OR}, \texttt{NOT}, \texttt{LT}, \texttt{GT}, \texttt{EQ}, etc.
    \item Control de flujo: saltos incondicionales (\texttt{JMP}) y condicionales (\texttt{JZ}, \texttt{JNZ}).
    \item Subrutinas y sistema: \texttt{CALL}, \texttt{RET}, \texttt{PRINT}, \texttt{READ}.
\end{itemize}

\section{Pruebas}
Para validar exhaustivamente el compilador, se diseñó una batería de pruebas en C (\texttt{test1.c} a \texttt{test5.c}) que evalúan el comportamiento del analizador ante diferentes niveles de estrés gramatical. Estas pruebas confirman el reconocimiento léxico, la correcta precedencia, la viabilidad de la tabla de símbolos y la precisión de la traducción intermedia y final.

\subsection{Pruebas realizadas}

\paragraph{Prueba 1: Precedencia matemática y operadores lógicos (\texttt{test1.c})}
Este archivo evalúa las expresiones aritméticas complejas, la precedencia forzada por paréntesis y la correcta agrupación de operadores booleanos \texttt{\&\&}, \texttt{||} y \texttt{!}. También incluye bucles para comprobar interrupciones condicionales.

\begin{lstlisting}[language=C, caption={Contenido de test1.c}]
int maxKeyframes = 100;
int currentFrame = 0;
float rotationSpeed = 2.5;

bool isShipActive = true;
bool hasCollided = false;
bool renderTorus = true;

// Prueba 1: Precedencia matematica estandar
// rotationSpeed * maxKeyframes debe evaluarse ANTES que la suma y resta
float newPosition = currentFrame + rotationSpeed * maxKeyframes - 10 / 2;

// Prueba 2: Precedencia forzada por parentesis
float offsetPosition = (currentFrame + rotationSpeed) * maxKeyframes;

print("Prueba de cadena en test1");
print("Iniciando simulacion de cuadros");

// Prueba 3: Operadores logicos, relacionales y NOT
// El AND (&&) debe agruparse antes que el OR (||)
// El NOT (!) debe resolverse de inmediato sobre renderTorus
if (currentFrame < maxKeyframes && isShipActive == true || !renderTorus) 
{
    currentFrame++;
} 
else if (hasCollided || currentFrame >= 10 && renderTorus) 
{
    isShipActive = false;
}

print("Estado de la nave actualizado");

// Prueba 4: Bucle while con multiples condiciones booleanas
while (isShipActive && !hasCollided) 
{
    currentFrame = currentFrame + 1;
    
    if (currentFrame == 50) 
    {
        hasCollided = true;
    }
}

print("Bucle principal finalizado");

// Prueba 5: For validando logica de interrupcion booleana
for (int i = 0; i < maxKeyframes; i++) 
{
    if (!isShipActive) 
    {
        break;
    }
}
\end{lstlisting}

\paragraph{Prueba 2: Funciones e I/O (\texttt{test2.c})}
Se enfoca en la declaración y llamada de funciones con y sin valor de retorno. Valida la captura de datos por teclado con \texttt{read}, la salida con \texttt{print} y la correcta interrupción del flujo con regresos explícitos.

\begin{lstlisting}[language=C, caption={Contenido de test2.c}]
int maxFrames = 60;
int currentFrame = 0;
bool engineReady = false;
int shipVelocity;

// Prueba 1: Funcion VOID con I/O y return vacio
void initGraphicsEngine()
{
    engineReady = true;
    print("Motor grafico inicializado");
    print(engineReady);
    return;
}

// Prueba 2: Funcion con tipo de retorno (INT), lectura de datos y return con operacion
int getShipSpeed()
{
    int inputSpeed;
    
    print("Esperando velocidad de entrada");
    read(inputSpeed); // Simulamos leer el input del teclado
    
    if (inputSpeed < 0)
    {
        print("Velocidad invalida detectada");
        return 0; // Return con valor estatico
    }
    
    print("Velocidad valida, aplicando multiplicador");
    return inputSpeed * 2; // Return con operacion matematica
}

// Prueba 3: Llamadas a funciones dentro de estatutos e interrupciones
void updateScene()
{
    if (!engineReady)
    {
        print("Motor no listo, se omite la escena");
        return; // Rompemos la ejecucion si el motor no esta listo
    }
    
    // Asignacion de una variable utilizando el valor de retorno de una funcion
    shipVelocity = getShipSpeed();
    
    currentFrame = currentFrame + 1;
    print("Escena actualizada correctamente");
    print(currentFrame);
}

// Prueba 4: Estructura de control llamando a funciones aisladas
void mainRenderLoop()
{
    initGraphicsEngine(); // Llamada a funcion como instruccion solitaria
    
    while (currentFrame < maxFrames && engineReady == true)
    {
        updateScene();
    }
}
\end{lstlisting}

\paragraph{Prueba 3: Ciclos anidados y optimización de código muerto (\texttt{test3.c})}
Este archivo concentra ciclos \texttt{do-while} y llamadas a funciones. Su propósito principal es generar escenarios de código muerto (instrucciones posteriores a un \texttt{return}) para verificar la limpieza de la salida intermedia durante la fase de optimización.

\begin{lstlisting}[language=C, caption={Contenido de test3.c}]
int helper()
{
    print("Calculando valor auxiliar");
    return 4;
}

int main()
{
    int a = 2;
    int b = 3;
    int c = a + b;
    int d = c * 4;
    bool flag = true;

    print("Inicio de test de expresiones");
    print(d);

    if (1 && 0 || !0)
    {
        d = d + 0;
        a = a * 1;
    }
    else
    {
        d = 99;
    }

    print("Bloque condicional evaluado");

    while (0)
    {
        a = a + 1;
        d = d + 2;
    };

    do
    {
        b = b + 1;
    }
    while (false);

    print("Bucle do-while completado");

    c = helper();
    print(c);

    print("Fin de test principal");

    return d;
    a = 100;
}
\end{lstlisting}

\paragraph{Prueba 4: Estructuras y selección múltiple (\texttt{test4.c})}
Representa un escenario de mayor complejidad gramatical. Incorpora la anidación de condicionales con declaraciones de estructuras \texttt{struct} y pone a prueba el mecanismo completo de bifurcación mediante \texttt{switch}, evaluando \texttt{case}, \texttt{default} y escapes con \texttt{break}.

\begin{lstlisting}[language=C, caption={Contenido de test4.c}]
struct EngineCore
{
    int throttle;
    int temperature;
    bool armed;
};

struct NavigationPack
{
    float altitude;
    float velocity;
    char mode;
    bool locked;
};

int boostLevel = 3;
int safetyMargin = 7;
bool telemetryReady = true;
char autopilotMode = 'A';

int calibrateCore()
{
    int base = 2;
    int factor = 5;
    int total = base + factor * 2;

    print("Calibrando nucleo principal");

    if (total > 10 && telemetryReady == true)
    {
        total = total - 1;
    }
    else
    {
        total = total + 1;
    }

    while (total < 20)
    {
        total = total + 2;
        if (total == 16)
        {
            break;
        }
    }

    print("Calibracion principal completada");

    return total;
}

void runDiagnostics()
{
    int iteration = 0;
    int status = 0;
    int coreResult = calibrateCore();

    print(coreResult);
    print("Ejecutando diagnosticos del sistema");

    for (int i = 0; i < 4; i++)
    {
        iteration = iteration + 1;
        status = status + i;

        if (i == 2)
        {
            status = status + 5;
        }
    }

    do
    {
        iteration = iteration - 1;
        status = status + 1;
    }
    while (iteration > 1);

    print("Diagnostico de ciclo completado");

    switch (status)
    {
        case 0:
            print(status);
            status = status + 1;
            break;
        case 3:
            if (status < 10)
            {
                status = status + 10;
            }
            else
            {
                status = status + 20;
            }
            break;
        default:
            while (status < 25)
            {
                status = status + 4;
            }
            break;
    }

    return;
}

int main()
{
    int initial = 1;
    int processed = 0;
    bool engineOn = true;

    print("Arrancando rutina principal");
    read(initial);

    processed = calibrateCore();
    if (initial < processed || !engineOn)
    {
        processed = processed + boostLevel;
    }
    else if (telemetryReady && engineOn)
    {
        processed = processed + safetyMargin;
    }
    else
    {
        processed = processed - 1;
    }

    runDiagnostics();
    print(processed);
    print("Rutina principal finalizada");

    return processed;
}
\end{lstlisting}

\paragraph{Prueba 5: Simulación de misión de vuelo (\texttt{test5.c})}
Es el archivo más riguroso, diseñado para empujar los límites del analizador. Integra la declaración de múltiples estructuras, funciones encadenadas que alteran el estado global, condicionales lógicos extendidos, bucles anidados y bloques de selección múltiple.

\begin{lstlisting}[language=C, caption={Contenido de test5.c}]
struct FlightControl
{
    int speed;
    int altitude;
    bool enabled;
};

struct PayloadBay
{
    int mass;
    char status;
    bool sealed;
};

int globalCounter = 0;
int targetAltitude = 120;
bool missionActive = true;
bool payloadReady = false;
char missionPhase = 'B';

int computeEnergy()
{
    int base = 8;
    int multiplier = 3;
    int energy = base * multiplier + 4;

    print("Calculando energia de vuelo");

    if (energy >= 28)
    {
        energy = energy - 2;
    }
    else
    {
        energy = energy + 2;
    }

    return energy;
}

void manageFlight()
{
    int cycle = 0;
    int flightEnergy = computeEnergy();

    print("Gestionando secuencia de vuelo");

    while (cycle < 3 && missionActive == true)
    {
        cycle = cycle + 1;
        globalCounter = globalCounter + cycle;

        if (cycle == 1)
        {
            payloadReady = true;
        }
        else if (cycle == 2)
        {
            payloadReady = false;
        }
        else
        {
            payloadReady = true;
        }
    }

    do
    {
        flightEnergy = flightEnergy - 1;
        targetAltitude = targetAltitude + 1;
    }
    while (flightEnergy > 20);

    print("Ajustando modo de mision");

    switch (missionPhase)
    {
        case 'A':
            print(globalCounter);
            break;
        case 'B':
            if (missionActive && payloadReady == false)
            {
                targetAltitude = targetAltitude + 5;
            }
            else
            {
                targetAltitude = targetAltitude + 10;
            }
            break;
        default:
            while (globalCounter < 3)
            {
                globalCounter = globalCounter + 1;
            }
            break;
    }

    return;
}

int missionControl()
{
    int result = 0;
    int localValue = 5;

    print("Iniciando control de mision");
    result = computeEnergy();
    result = result + localValue;

    if (result > targetAltitude && payloadReady == false || missionActive)
    {
        result = result + 1;
    }

    for (int step = 0; step < 5; step++)
    {
        if (step == 4)
        {
            break;
        }

        result = result + step;
    }

    manageFlight();
    print(result);
    print("Control de mision completado");

    return result;
}
\end{lstlisting}

\subsection{Archivos generados por el compilador}
Durante la compilación de cualquiera de los archivos de prueba anteriores, el sistema exporta de manera automática los siguientes reportes dentro de la carpeta \texttt{output/}:
\begin{itemize}[label={-}]
    \item \texttt{tabla\_tokens.txt}: Reporte del lexer con la numeración, lexema y clase de cada token reconocido.
    \item \texttt{tabla\_variables.txt}: Resumen de clasificación léxica separando variables, cadenas y errores detectados.
    \item \texttt{tabla\_simbolos.txt}: Registro oficial del parser listando nombres, tipos (\texttt{int}, \texttt{struct}, \texttt{void}, etc.) y categorías.
    \item \texttt{codigo\_intermedio.txt}: El código fuente linealizado en formato de tres direcciones nativo.
    \item \texttt{codigo\_intermedio\_optimizado.txt}: La versión depurada del archivo anterior, resultado del plegado y la eliminación de redundancias.
    \item \texttt{codigo\_final.txt}: La traducción definitiva a los mnemónicos de ensamblador virtual.
\end{itemize}

El aislamiento de estas salidas en archivos independientes facilitó enormemente la auditoría y depuración manual durante el ciclo de vida del proyecto.

\section{Resultados obtenidos}
Los resultados de la ejecución sobre la batería de pruebas confirman que el compilador cumple satisfactoriamente con los objetivos planteados:
\begin{itemize}[label={-}]
    \item El lenguaje extendido (con \texttt{struct} y \texttt{switch}) es analizado léxica y sintácticamente sin conflictos o ambigüedades.
    \item Se extrae y consolida información semántica precisa en la tabla de símbolos.
    \item El proceso de traducción produce código intermedio lógico, secuencial y completamente funcional.
    \item El módulo de optimización demuestra la capacidad de simplificar expresiones aritméticas estáticas y reducir la entropía del código generado.
    \item El seudoensamblador final refleja un mapeo fiel a la lógica concebida en el programa original en C.
\end{itemize}

\section{Conclusiones}
Este proyecto logró integrar con éxito todas las fases teóricas del proceso de compilación dentro de una herramienta funcional y cohesionada. Aunque la implementación mantiene un carácter académico, la cobertura de los lineamientos del curso es total: desde el reconocimiento de expresiones regulares hasta la emisión de código final, pasando por la indispensable traducción dirigida por la sintaxis.

El desarrollo iterativo dejó como gran aprendizaje que el reto de un compilador no radica únicamente en escribir reglas de gramática correctas, sino en orquestar el flujo de datos: gestionar lexemas, resolver precedencias lógicas cruzadas y asegurar que los atributos se propaguen sin pérdida a través del árbol de derivación.

Como prospectiva técnica, las mejoras naturales para iteraciones futuras incluirían la implementación del paso de parámetros dinámicos hacia funciones, el soporte nativo para algoritmos recursivos, la verificación estricta del sistema de tipos (Type Checking) previo a la generación de código, y el desarrollo de una máquina virtual (VM) en C capaz de ingerir y ejecutar directamente el archivo \texttt{codigo\_final.txt}.

\section{Tecnologías utilizadas}
\begin{itemize}[label={-}]
    \item \textbf{C}: Lenguaje anfitrión para la lógica del compilador y estructuras de datos subyacentes.
    \item \textbf{Flex}: Herramienta de generación del analizador léxico.
    \item \textbf{Bison}: Herramienta de generación del analizador sintáctico (LALR).
    \item \textbf{LaTeX}: Sistema de composición de textos empleado para la documentación técnica.
\end{itemize}

\end{document}